% Endpoint projection: one global choreography is compiled into one local
% program per role, adapted from Wiersdorf's ChoRex tutorial.
% Requires: \usepackage{tikz}
%   \usetikzlibrary{arrows.meta,positioning,fit,backgrounds}
% Include with a normal figure environment, e.g. % Endpoint projection: one global choreography is compiled into one local
% program per role, adapted from Wiersdorf's ChoRex tutorial.
% Requires: \usepackage{tikz}
%   \usetikzlibrary{arrows.meta,positioning,fit,backgrounds}
% Include with a normal figure environment, e.g. % Endpoint projection: one global choreography is compiled into one local
% program per role, adapted from Wiersdorf's ChoRex tutorial.
% Requires: \usepackage{tikz}
%   \usetikzlibrary{arrows.meta,positioning,fit,backgrounds}
% Include with a normal figure environment, e.g. % Endpoint projection: one global choreography is compiled into one local
% program per role, adapted from Wiersdorf's ChoRex tutorial.
% Requires: \usepackage{tikz}
%   \usetikzlibrary{arrows.meta,positioning,fit,backgrounds}
% Include with a normal figure environment, e.g. \input{figures/epp}.
\begin{tikzpicture}[
  >=Latex, font=\small,
  code/.style={align=left,font=\scriptsize\ttfamily,inner sep=0pt},
  title/.style={font=\small\bfseries,inner sep=0pt},
  prog/.style={rounded corners,draw,fill=gray!4,inner sep=5pt},
  local/.style={prog,fill=gray!12},
  split/.style={->,thick,gray!70},
]

% ---------- global choreography (top) ----------
\node[code] (ccode) at (-0.55,4.95) {%
  Alice.get\_modulus() \textasciitilde> Bob.(m)\\
  Bob.make\_session(m) \textasciitilde> Alice.(session)\\
  Alice.encrypt(session, message) \textasciitilde> Bob.(c)\\
  Bob.decrypt(c)};
\node[title,anchor=south] (ctitle) at ([yshift=2.5mm]ccode.north) {Choreography};
\begin{scope}[on background layer]
  \node[prog,fit=(ctitle)(ccode)] (chor) {};
\end{scope}

% ---------- projected local programs (bottom) ----------
\node[code,anchor=north east] (acode) at (-0.65,1.15) {%
  send(Bob, get\_modulus())\\
  session = receive()\\
  send(Bob, encrypt(session, message))};
\node[title,anchor=south] (atitle) at ([yshift=2.5mm]acode.north) {Code for Alice};
\begin{scope}[on background layer]
  \node[local,fit=(atitle)(acode)] (alice) {};
\end{scope}

\node[code,anchor=north west] (bcode) at (0.65,1.15) {%
  m = receive()\\
  send(Alice, make\_session(m))\\
  c = receive()\\
  decrypt(c)};
\node[title,anchor=south] (btitle) at ([yshift=2.5mm]bcode.north) {Code for Bob};
\begin{scope}[on background layer]
  \node[local,fit=(btitle)(bcode)] (bob) {};
\end{scope}

% ---------- projection ----------
\draw[split] ([xshift=-6mm]chor.south) -- (alice.north);
\draw[split] ([xshift= 6mm]chor.south) -- (bob.north);
\node[align=center,font=\small\bfseries] at (-0.55,3.0) {Endpoint\\projection};

\end{tikzpicture}
.
\begin{tikzpicture}[
  >=Latex, font=\small,
  code/.style={align=left,font=\scriptsize\ttfamily,inner sep=0pt},
  title/.style={font=\small\bfseries,inner sep=0pt},
  prog/.style={rounded corners,draw,fill=gray!4,inner sep=5pt},
  local/.style={prog,fill=gray!12},
  split/.style={->,thick,gray!70},
]

% ---------- global choreography (top) ----------
\node[code] (ccode) at (-0.55,4.95) {%
  Alice.get\_modulus() \textasciitilde> Bob.(m)\\
  Bob.make\_session(m) \textasciitilde> Alice.(session)\\
  Alice.encrypt(session, message) \textasciitilde> Bob.(c)\\
  Bob.decrypt(c)};
\node[title,anchor=south] (ctitle) at ([yshift=2.5mm]ccode.north) {Choreography};
\begin{scope}[on background layer]
  \node[prog,fit=(ctitle)(ccode)] (chor) {};
\end{scope}

% ---------- projected local programs (bottom) ----------
\node[code,anchor=north east] (acode) at (-0.65,1.15) {%
  send(Bob, get\_modulus())\\
  session = receive()\\
  send(Bob, encrypt(session, message))};
\node[title,anchor=south] (atitle) at ([yshift=2.5mm]acode.north) {Code for Alice};
\begin{scope}[on background layer]
  \node[local,fit=(atitle)(acode)] (alice) {};
\end{scope}

\node[code,anchor=north west] (bcode) at (0.65,1.15) {%
  m = receive()\\
  send(Alice, make\_session(m))\\
  c = receive()\\
  decrypt(c)};
\node[title,anchor=south] (btitle) at ([yshift=2.5mm]bcode.north) {Code for Bob};
\begin{scope}[on background layer]
  \node[local,fit=(btitle)(bcode)] (bob) {};
\end{scope}

% ---------- projection ----------
\draw[split] ([xshift=-6mm]chor.south) -- (alice.north);
\draw[split] ([xshift= 6mm]chor.south) -- (bob.north);
\node[align=center,font=\small\bfseries] at (-0.55,3.0) {Endpoint\\projection};

\end{tikzpicture}
.
\begin{tikzpicture}[
  >=Latex, font=\small,
  code/.style={align=left,font=\scriptsize\ttfamily,inner sep=0pt},
  title/.style={font=\small\bfseries,inner sep=0pt},
  prog/.style={rounded corners,draw,fill=gray!4,inner sep=5pt},
  local/.style={prog,fill=gray!12},
  split/.style={->,thick,gray!70},
]

% ---------- global choreography (top) ----------
\node[code] (ccode) at (-0.55,4.95) {%
  Alice.get\_modulus() \textasciitilde> Bob.(m)\\
  Bob.make\_session(m) \textasciitilde> Alice.(session)\\
  Alice.encrypt(session, message) \textasciitilde> Bob.(c)\\
  Bob.decrypt(c)};
\node[title,anchor=south] (ctitle) at ([yshift=2.5mm]ccode.north) {Choreography};
\begin{scope}[on background layer]
  \node[prog,fit=(ctitle)(ccode)] (chor) {};
\end{scope}

% ---------- projected local programs (bottom) ----------
\node[code,anchor=north east] (acode) at (-0.65,1.15) {%
  send(Bob, get\_modulus())\\
  session = receive()\\
  send(Bob, encrypt(session, message))};
\node[title,anchor=south] (atitle) at ([yshift=2.5mm]acode.north) {Code for Alice};
\begin{scope}[on background layer]
  \node[local,fit=(atitle)(acode)] (alice) {};
\end{scope}

\node[code,anchor=north west] (bcode) at (0.65,1.15) {%
  m = receive()\\
  send(Alice, make\_session(m))\\
  c = receive()\\
  decrypt(c)};
\node[title,anchor=south] (btitle) at ([yshift=2.5mm]bcode.north) {Code for Bob};
\begin{scope}[on background layer]
  \node[local,fit=(btitle)(bcode)] (bob) {};
\end{scope}

% ---------- projection ----------
\draw[split] ([xshift=-6mm]chor.south) -- (alice.north);
\draw[split] ([xshift= 6mm]chor.south) -- (bob.north);
\node[align=center,font=\small\bfseries] at (-0.55,3.0) {Endpoint\\projection};

\end{tikzpicture}
.
\begin{tikzpicture}[
  >=Latex, font=\small,
  code/.style={align=left,font=\scriptsize\ttfamily,inner sep=0pt},
  title/.style={font=\small\bfseries,inner sep=0pt},
  prog/.style={rounded corners,draw,fill=gray!4,inner sep=5pt},
  local/.style={prog,fill=gray!12},
  split/.style={->,thick,gray!70},
]

% ---------- global choreography (top) ----------
\node[code] (ccode) at (-0.55,4.95) {%
  Alice.get\_modulus() \textasciitilde> Bob.(m)\\
  Bob.make\_session(m) \textasciitilde> Alice.(session)\\
  Alice.encrypt(session, message) \textasciitilde> Bob.(c)\\
  Bob.decrypt(c)};
\node[title,anchor=south] (ctitle) at ([yshift=2.5mm]ccode.north) {Choreography};
\begin{scope}[on background layer]
  \node[prog,fit=(ctitle)(ccode)] (chor) {};
\end{scope}

% ---------- projected local programs (bottom) ----------
\node[code,anchor=north east] (acode) at (-0.65,1.15) {%
  send(Bob, get\_modulus())\\
  session = receive()\\
  send(Bob, encrypt(session, message))};
\node[title,anchor=south] (atitle) at ([yshift=2.5mm]acode.north) {Code for Alice};
\begin{scope}[on background layer]
  \node[local,fit=(atitle)(acode)] (alice) {};
\end{scope}

\node[code,anchor=north west] (bcode) at (0.65,1.15) {%
  m = receive()\\
  send(Alice, make\_session(m))\\
  c = receive()\\
  decrypt(c)};
\node[title,anchor=south] (btitle) at ([yshift=2.5mm]bcode.north) {Code for Bob};
\begin{scope}[on background layer]
  \node[local,fit=(btitle)(bcode)] (bob) {};
\end{scope}

% ---------- projection ----------
\draw[split] ([xshift=-6mm]chor.south) -- (alice.north);
\draw[split] ([xshift= 6mm]chor.south) -- (bob.north);
\node[align=center,font=\small\bfseries] at (-0.55,3.0) {Endpoint\\projection};

\end{tikzpicture}
